\chapter[Referencial Teórico]{Referencial Teórico}

\section{Trabalho de Conclusão de Curso}

O Trabalho de Conclusão de Curso desempenha um papel
importante na formação acadêmica, proporcionando ao
estudante a oportunidade de desenvolver atividades de
pesquisa e aplicar conhecimentos adquiridos durante a
graduação. Segundo \citeonline{Trindade2018}, o TCC
pode contribuir para a formação de profissionais com
maior capacidade crítica e reflexiva, ao estimular a
investigação e a construção do conhecimento.

Além da produção acadêmica, o desenvolvimento do TCC
está relacionado à autonomia do estudante.
\citeonline{Trindade2018} destacam que a realização
de atividades de pesquisa favorece o desenvolvimento
da iniciativa e da capacidade de analisar questões
de maneira crítica, competências que também podem
contribuir para a atuação profissional.

Os autores também ressaltam que o TCC não precisa
necessariamente apresentar uma descoberta inédita,
podendo abordar assuntos relacionados à área de
formação do estudante. Nesse sentido, sua relevância
não está apenas no resultado apresentado, mas também
nas experiências de aprendizagem proporcionadas
durante sua elaboração.


\subsection{Organização do TCC na FCTE}

Na Faculdade de Ciências e Tecnologias em Engenharia (FCTE)
da Universidade de Brasília (UnB), o Trabalho de Conclusão
de Curso é regulamentado pela Resolução do Conselho da FGA
nº 0002/2022, que estabelece sua organização em dois
componentes curriculares: TCC 1 e TCC 2 \cite{UnB2022}.

O TCC 1 possui caráter qualificatório, tendo como objetivo
avaliar a aptidão e a maturidade do discente e do tema
escolhido em sua área de investigação. Já o TCC 2 corresponde à versão definitiva do trabalho, destinada à sua apresentação e validação. Ambos os componentes são avaliados por bancas examinadoras específicas \cite{UnB2022}.

Após a avaliação, o resultado é registrado em ata e, no caso do TCC 2 aprovado, a versão definitiva é encaminhada para disponibilização na Biblioteca Digital de Monografias (BDM) \cite{UnB2022}.

Além das atividades realizadas pelos estudantes e orientadores, a coordenação também participa dos procedimentos acadêmicos
relacionados ao TCC, incluindo a aprovação da composição das bancas, o acompanhamento dos trabalhos, o controle dos prazos e a reunião das atas e menções ao final de cada semestre \cite{UnB2022}.

\section{Sistemas de Informação no Ensino Superior}

\subsection{Sistemas de Informação na Gestão Acadêmica}

\section{Trabalhos e Soluções Relacionadas}

\section{Análise e Discussão}